\documentclass[UTF8,11pt]{ctexart}
\usepackage[a4paper,margin=24mm]{geometry}
\usepackage{amsmath,amssymb,amsthm,mathtools,mathrsfs}
\usepackage[all]{xy}
\usepackage{xurl,hyperref,enumitem}
\usepackage{tikz}
\usetikzlibrary{arrows.meta,cd,decorations.markings,calc,patterns}
\hypersetup{hidelinks,breaklinks=true}
\setlength{\emergencystretch}{3em}
\allowdisplaybreaks
\newtheorem*{theorem}{Theorem}
\newtheorem*{thm}{Theorem}
\newtheorem*{lemma}{Lemma}
\newtheorem*{proposition}{Proposition}
\newtheorem*{prop}{Proposition}
\newtheorem*{corollary}{Corollary}
\newtheorem*{cor}{Corollary}
\newtheorem*{claim}{Claim}
\theoremstyle{definition}
\newtheorem*{definition}{Definition}
\newtheorem*{defn}{Definition}
\newtheorem*{remark}{Remark}
\newtheorem*{example}{Example}
\newcommand{\R}{\mathbb R}
\newcommand{\C}{\mathbb C}
\newcommand{\Z}{\mathbb Z}
\newcommand{\Q}{\mathbb Q}
\newcommand{\N}{\mathbb N}
\newcommand{\F}{\mathbb F}
\newcommand{\D}{\mathbb D}
\newcommand{\bB}{\mathbb B}
\newcommand{\bC}{\mathbb C}
\newcommand{\bR}{\mathbb R}
\newcommand{\bZ}{\mathbb Z}
\newcommand{\bN}{\mathbb N}
\newcommand{\bQ}{\mathbb Q}
\newcommand{\bD}{\mathbb D}
\newcommand{\bF}{\mathbb F}
\newcommand{\CP}{\mathbb{CP}}
\newcommand{\RP}{\mathbb{RP}}
\newcommand{\sO}{\mathscr O}
\newcommand{\sI}{\mathscr I}
\newcommand{\sF}{\mathscr F}
\newcommand{\sG}{\mathscr G}
\newcommand{\sH}{\mathscr H}
\newcommand{\sR}{\mathscr R}
\newcommand{\sM}{\mathscr M}
\newcommand{\sV}{\mathscr V}
\newcommand{\sU}{\mathscr U}
\DeclareMathOperator{\Hom}{Hom}
\DeclareMathOperator{\Ext}{Ext}
\DeclareMathOperator{\Tor}{Tor}
\DeclareMathOperator{\Spec}{Spec}
\DeclareMathOperator{\Proj}{Proj}
\DeclareMathOperator{\supp}{supp}
\DeclareMathOperator{\im}{im}
\DeclareMathOperator{\id}{id}
\DeclareMathOperator{\Tr}{Tr}
\DeclareMathOperator{\tr}{tr}
\DeclareMathOperator{\rank}{rank}
\DeclareMathOperator{\ord}{ord}
\DeclareMathOperator{\dist}{dist}
\DeclareMathOperator{\End}{End}
\DeclareMathOperator{\Aut}{Aut}
\DeclareMathOperator{\Gal}{Gal}
\DeclareMathOperator{\vol}{vol}
\DeclareMathOperator{\Cl}{Cl}
\DeclareMathOperator{\coker}{coker}
\DeclareMathOperator{\codim}{codim}
\DeclareMathOperator{\divergence}{div}
\DeclareMathOperator{\Ric}{Ric}
\DeclareMathOperator{\grad}{grad}
\DeclareMathOperator{\Hess}{Hess}
\newcommand{\abs}[1]{\left\lvert#1\right\rvert}
\newcommand{\sabs}[1]{\lvert#1\rvert}
\newcommand{\babs}[1]{\bigl\lvert#1\bigr\rvert}
\newcommand{\Babs}[1]{\Bigl\lvert#1\Bigr\rvert}
\newcommand{\bbabs}[1]{\biggl\lvert#1\biggr\rvert}
\newcommand{\BBabs}[1]{\Biggl\lvert#1\Biggr\rvert}
\newcommand{\norm}[1]{\left\lVert#1\right\rVert}
\newcommand{\snorm}[1]{\lVert#1\rVert}
\newcommand{\bnorm}[1]{\bigl\lVert#1\bigr\rVert}
\newcommand{\Bnorm}[1]{\Bigl\lVert#1\Bigr\rVert}
\newcommand{\bbnorm}[1]{\biggl\lVert#1\biggr\rVert}
\newcommand{\BBnorm}[1]{\Biggl\lVert#1\Biggr\rVert}
\newcommand{\widebar}[1]{\overline{#1}}
\newcommand{\nicefrac}[2]{\frac{#1}{#2}}
\newcommand{\linnprod}[2]{\langle#1,#2\rangle}
\newcommand{\myindex}[1]{#1}
\newcommand{\glsadd}[1]{}
\newcommand{\avoidbreak}{}
\newcommand{\sourceref}[1]{\textnormal{[原文：\nolinkurl{#1}]}}
\newcommand{\thmref}[1]{\sourceref{#1}}
\newcommand{\lemmaref}[1]{\sourceref{#1}}
\newcommand{\propref}[1]{\sourceref{#1}}
\newcommand{\corref}[1]{\sourceref{#1}}
\newcommand{\defnref}[1]{\sourceref{#1}}
\newcommand{\exampleref}[1]{\sourceref{#1}}
\newcommand{\exerciseref}[1]{\sourceref{#1}}
\newcommand{\figureref}[1]{\sourceref{#1}}
\newcommand{\sectionref}[1]{\sourceref{#1}}
\newcommand{\appendixref}[1]{\sourceref{#1}}
\newcommand{\stacksref}[2]{\href{https://stacks.math.columbia.edu/tag/#1}{#2 [#1]}}


\newcommand{\E}{\mathbb E}
\newcommand{\Prob}{\mathbb P}
\newcommand{\Reff}{\mathcal R_{\mathrm{eff}}}
\newcommand{\eps}{\varepsilon}
\DeclareMathOperator{\diam}{diam}
\DeclareMathOperator{\Var}{Var}
\DeclareMathOperator{\Crit}{Crit}
\newcommand{\dd}{\,\mathrm d}

\begin{document}
\section*{058\quad 双圆盘到单位球不存在适当全纯映射}
\subsection*{来源与计数目标}
Ji\v{r}\'i Lebl, \emph{Tasty Bits of Several Complex Variables}, Chapter 2, ``Inequivalence of ball and polydisc''。原生 \texttt{scv.tex}：\texttt{thm:Rothstein}、\texttt{prop:noanaldiscinsphere}、\texttt{lemma:bndrytobndry} 及主证明。使用 commit \texttt{1149c5827e3ee1e25cd0939bfa846ca441517e62}，读取范围包含第3100--3370行；获取日期2026-09-28。
\par\url{https://github.com/jirilebl/scv/blob/1149c5827e3ee1e25cd0939bfa846ca441517e62/scv.tex}
\par\textbf{本文件只计一个问题：}证明不存在 $\D^2\to\bB_2$ 的适当全纯映射；不另计球面无解析圆盘和适当映射边界判据。
\subsection*{作者命题与人类证明}
\begin{definition}[Proper map]
Suppose $f\colon X\to Y$ is a continuous map between two topological spaces. Then $f$ is a \emph{proper map} if for every compact $K\subset\subset Y$, the set $f^{-1}(K)$ is compact.
\end{definition}
\begin{theorem}[Rothstein 1935]\label{thm:Rothstein}
There exists no proper holomorphic mapping of the unit bidisc $\D^2=\D\times\D\subset\C^2$ to the unit ball $\bB_2\subset\C^2$.
\end{theorem}
\begin{definition}[Analytic disc]
A nonconstant holomorphic mapping $\varphi\colon\D\to\C^n$ is called an \emph{analytic disc}. If the mapping $\varphi$ extends continuously to the closed unit disc $\widebar\D$, then the mapping $\varphi\colon\widebar\D\to\C^n$ is called a \emph{closed analytic disc}.
Often we call the image $\Delta=\varphi(\D)$ the analytic disc rather than the mapping. For a closed analytic disc we write $\partial\Delta=\varphi(\partial\D)$ and call it the boundary of the analytic disc.
\end{definition}
\begin{proposition}\label{prop:noanaldiscinsphere}
The unit sphere $S^{2n-1}=\partial\bB_n\subset\C^n$ contains no analytic discs.
\end{proposition}
\begin{proof}
Suppose there is a holomorphic function $g\colon\D\to\C^n$ such that the image $g(\D)$ is contained in the unit sphere. In other words, for all $z\in\D$,
\[\snorm{g(z)}^2=\sabs{g_1(z)}^2+\sabs{g_2(z)}^2+\cdots+\sabs{g_n(z)}^2=1.\]
Without loss of generality (after composing with a unitary matrix), assume that $g(0)=(1,0,0,\ldots,0)$. Consider the first component and notice that $g_1(0)=1$. If a sum of nonnegative numbers is less than or equal to $1$, then they all are. Hence, $\sabs{g_1(z)}\leq1$. The maximum principle says that $g_1(z)=1$ for all $z\in\D$. But then $g_k(z)=0$ for all $k=2,\ldots,n$ and all $z\in\D$. Therefore, $g$ is constant and thus not an analytic disc.
\end{proof}
Before we prove the theorem, let us make the statement about proper maps taking boundary to boundary precise.
\begin{lemma}\label{lemma:bndrytobndry}
Let $U\subset\R^n$ and $V\subset\R^m$ be bounded domains and let $f\colon U\to V$ be continuous. Then $f$ is proper if and only if for every sequence $\{p_k\}$ in $U$ such that $p_k\to p\in\partial U$, the set of limit points of $\bigl\{f(p_k)\bigr\}$ lies in $\partial V$.
\end{lemma}
\begin{proof}
Suppose $f$ is proper. Let $\{p_k\}$ be a sequence in $U$ such that $p_k\to p\in\partial U$. Take any convergent subsequence $\bigl\{f(p_{k_\ell})\bigr\}$ of $\bigl\{f(p_k)\bigr\}$ converging to some $q\in\widebar V$. Consider $E=\bigl\{f(p_{k_\ell})\bigr\}$ as a set. Let $\widebar E$ be the closure of $E$ in $V$ (subspace topology). The inverse image $f^{-1}(\widebar E)$ is not compact (it contains a sequence going to $p\in\partial U$) and hence $\widebar E=(E\cup\{q\})\cap V$ is not compact either as $f$ is proper. That means that $q\notin V$ and so $q\in\partial V$. As we took an arbitrary convergent subsequence of $\bigl\{f(p_k)\bigr\}$, $q$ was an arbitrary limit point. Therefore, all limit points are in $\partial V$.

Let us prove the converse. Suppose that for every sequence $\{p_k\}$ in $U$ such that $p_k\to p\in\partial U$, the set of limit points of $\bigl\{f(p_k)\bigr\}$ lies in $\partial V$. Take a closed set $E\subset V$ (subspace topology) and suppose $f^{-1}(E)$ is not compact. Then there exists a sequence $\{p_k\}$ in $f^{-1}(E)$ such that $p_k\to p\in\partial U$, because $f^{-1}(E)$ is closed (in $U$), bounded, but not compact. The hypothesis then says that the limit points of $\bigl\{f(p_k)\bigr\}$ are in $\partial V$. Hence, $E$ has limit points in $\partial V$ and is not compact.
\end{proof}
\begin{proof}[Proof of Rothstein's theorem]
Suppose there is a proper holomorphic map $f\colon\D^2\to\bB_2$. Fix some $e^{i\theta}$ in the boundary of the disc $\D$. Take a sequence $w_k\in\D$ such that $w_k\to e^{i\theta}$. The functions $g_k(\zeta)=f(\zeta,w_k)$ map the unit disc into $\bB_2$. By Montel's theorem and by passing to a subsequence, assume that the sequence of functions converges (uniformly on compact subsets) to a limit $g\colon\D\to\overline{\bB}_2$. As $(\zeta,w_k)\to(\zeta,e^{i\theta})\in\partial\D^2$, by Lemma~\ref{lemma:bndrytobndry}, $g(\D)\subset\partial\bB_2$, and hence $g$ is constant by Proposition~\ref{prop:noanaldiscinsphere}.

Let $g_k'$ denote the derivative (we differentiate each component). The functions $g_k'$ converge to $g'=0$. So for an arbitrary fixed $\zeta\in\D$, $\frac{\partial f}{\partial z_1}(\zeta,w_k)\to0$. This limit holds for all $e^{i\theta}$ and some subsequence of an arbitrary sequence $\{w_k\}$ where $w_k\to e^{i\theta}$. The holomorphic mapping $w\mapsto\frac{\partial f}{\partial z_1}(\zeta,w)$, therefore, extends continuously to the closure $\widebar\D$ and is zero on $\partial\D$. We apply the maximum principle or the Cauchy formula and the fact that $\zeta$ was arbitrary to find $\frac{\partial f}{\partial z_1}\equiv0$. By symmetry $\frac{\partial f}{\partial z_2}\equiv0$. Therefore, $f$ is constant, which is a contradiction as $f$ was proper.
\end{proof}
\subsection*{整理说明（非作者正文）}
作者原命题与全部推导保留；删去无关练习、历史讨论，以及证明结束后重复解释同一切片构造的示意图段。图并不承担额外推导，边界切片、极限和导数论证全文在上。标准先修为 Montel 定理、紧集上一致收敛的全纯函数导数收敛、最大值原理及欧氏空间紧性。实际难点是选取趋向边界的解析切片，把适当性与球面几何转成偏导数的全局消失；不是直接引用球与双圆盘不等价。与033的圆域固定点线性化目标及证明不同。保留的是有定位的原生 TeX 摘录，不是整书下载件。
\subsection*{可修改的简短标注（非作者正文）}
$c$：多复变区域间适当全纯映射的不存在性

$a$：解析圆盘；适当映射与边界序列；正规族与Montel定理；切片极限；最大值原理

\texttt{cross\_theory\_status = yes}。将双圆盘的边界切片通过适当性送到球面，以正规族极限和球面无解析圆盘性质迫使切片导数为零，再返回全局映射矛盾。位置：Rothstein 主证明及两项所附支撑论证。
标签为非穷尽、可修改初稿。
\subsection*{许可与修改记录}
原作者及作品见来源栏；来源采用 CC BY-SA 4.0。本题证摘录和附加整理沿用该许可。2026-09-28：增加题号、来源定位、独立排版、整理说明和初稿标注；数学内容的取材范围及排版变化见整理说明。
\end{document}
